\subsection{Sampling Interval ($\Delta$t)}
Consistent sampling intervals are important factors in time series data. Large time jumps between consecutive individual measurements can introduce artificial anomalies. Therefore the dataset was validated by calculating the time differences between consecutive timestamps. 
The interval $\Delta t$ was calculated as : \\
\begin{center}
$\Delta t_n = t_n - t_{n-1}$ \\
\end{center}   
Where $\Delta t_n$ is the time  difference between the current  and previous measurement. 
With Python's NumPy min and max functions, the minimum and maximum $\Delta t$ were computed for each dataset. The complete evaluation results are detailed in \ref{app:appendix2}. Inspection identified anomalies characterized by unexpected negative or excessively large values. After excluding these specific files  
(Files: 05, 13, 14, and 16) and running again the algorithm, the resulting statistical values improved significantly, as show in \ref{table:statistics}.\\
\begin{table}[!h]
 \centering
 \begin{tabular}{|c|c|}
 \hline
 Total number of intervals & 2,975,473 \\ 
 \hline
 Mean $\Delta t$ & 1,906.94 $\mu$s \\ 
 \hline
 Standard deviation & 38.03 $\mu$s \\ 
 \hline
 Median & 1,912.0 $\mu$s \\
 \hline
 Minimum & 1,660.0 $\mu$s \\
 \hline
 Maximum & 2,032.0 $\mu$s \\
 \hline
 \end{tabular}
 \caption{Statistics of sampling intervals after filtering rounded to the second decimal}
 \label{table:statistics}
\end{table}
\\
By analyzing the results, a clear improvement can be observed, the difference between the minimum, maximum, and the global mean are very small. Additionally, the standard deviation has improved a lot, which indicates highly consistent sampling. These results imply that the dataset was successful cleaned and can be used for further benchmarking and testing.